\chapter{Bạn Bè Ở Tuổi Đi Học}
\markboth{Bạn Bè Ở Tuổi Đi Học}{Friendship During School Age}

\section{Một Người Bạn Đỡ Cả Một Năm Học}
\begin{truyen}{Một Người Bạn Đỡ Cả Một Năm Học}{One Friend Can Hold Up an Entire School Year}
\chuhoa{C}{ó những em đi học được không phải vì môn học dễ, mà vì trong lớp có một người bạn chịu ngồi cạnh.}
\textit{(Some students manage school not because classes are easy, but because there is one friend willing to sit beside them.)}

Một người rủ cùng ăn, cùng làm bài, cùng đi xuống phòng thi đôi khi đỡ cả tinh thần.
\textit{(One person who shares lunch, homework, and the walk to an exam can steady an entire spirit.)}

Người lớn thường đánh giá thấp sức đỡ của tình bạn tuổi học trò.
\textit{(Adults often underestimate how much school-age friendship can support a child.)}

Có khi một người bạn tốt cứu được thứ mà bài giảng không cứu được.
\textit{(Sometimes a good friend rescues what a lesson cannot.)}

\end{truyen}

\begin{baihoc}
Tình bạn đúng lúc cũng là một dạng chống đỡ tinh thần rất thật.
\textit{(Timely friendship is a very real form of emotional support.)}
\end{baihoc}

\begin{ghinhoanh}
\textbf{Short English to remember:}

Timely friendship is a very real form of emotional support.

\end{ghinhoanh}

\begin{tuhocnhanh}
\textbf{Gợi ý để dạy và tự nghĩ / Teaching and reflection prompts}

1. Một người bạn tốt đang đỡ điều gì? \textit{What can one good friend help carry?}

2. Vì sao người lớn hay đánh giá thấp điều này? \textit{Why do adults underestimate this?}

3. Lớp học có đang tạo cơ hội cho tình bạn tốt không? \textit{Is the classroom making room for healthy friendship?}
\end{tuhocnhanh}
\ngancach

\section{Bị Lôi Vào Một Nhóm Xấu}
\begin{truyen}{Bị Lôi Vào Một Nhóm Xấu}{Being Pulled into the Wrong Group}
\chuhoa{N}{hiều em ban đầu không xấu. Chỉ là quá sợ bị lạc nên đi theo một nhóm cho đỡ cô đơn.}
\textit{(Many students are not bad at the beginning. They are simply afraid of being alone and go with the group that will take them.)}

Muốn được thuộc về là một nhu cầu rất mạnh ở tuổi học trò.
\textit{(The need to belong is very strong during school years.)}

Khi chưa có chỗ đứng lành mạnh, các em dễ chọn chỗ đứng ồn ào.
\textit{(Without a healthy place to stand, children often choose a louder one.)}

Không ít hành vi sai bắt đầu từ nỗi sợ bị bỏ rơi.
\textit{(Many wrong behaviors begin with the fear of being left out.)}

\end{truyen}

\begin{baihoc}
Muốn giữ học sinh khỏi nhóm xấu, phải cho các em một chỗ thuộc về tốt hơn.
\textit{(If we want to keep students away from harmful groups, we must offer them a better place to belong.)}
\end{baihoc}

\begin{ghinhoanh}
\textbf{Short English to remember:}

Students need a better place to belong.

\end{ghinhoanh}

\begin{tuhocnhanh}
\textbf{Gợi ý để dạy và tự nghĩ / Teaching and reflection prompts}

1. Nỗi sợ nào đang kéo em này đi? \textit{What fear is pulling this student?}

2. Vì sao nhu cầu thuộc về lại mạnh đến vậy? \textit{Why is belonging so powerful at this age?}

3. Lớp học có thể tạo chỗ đứng lành mạnh ra sao? \textit{How can the classroom create healthier belonging?}
\end{tuhocnhanh}
\ngancach

\section{Mượn Tiền Rồi Mất Cả Bạn}
\begin{truyen}{Mượn Tiền Rồi Mất Cả Bạn}{Borrowed Money Can Cost a Friendship}
\chuhoa{C}{ó những tình bạn học trò tưởng rất bền, nhưng chỉ cần đụng đến tiền là lộ ra một tầng khác.}
\textit{(Some school friendships look strong until money enters the picture.)}

Một khoản nhỏ cũng đủ làm mất lòng, mất tin, rồi mất luôn cả sự tự nhiên cũ.
\textit{(Even a small amount can damage trust and the old ease between two students.)}

Học sinh còn nhỏ nên chưa biết nói rõ, chưa biết trả đúng, chưa biết giữ cam kết.
\textit{(Students are still young and may not yet know how to communicate, repay, or honor promises well.)}

Có những bài học về ranh giới đến rất sớm qua một món nợ nhỏ.
\textit{(Some lessons about boundaries arrive early through a small debt.)}

\end{truyen}

\begin{baihoc}
Tình bạn muốn bền cũng cần được đi cùng trách nhiệm.
\textit{(Friendship lasts better when responsibility walks with it.)}
\end{baihoc}

\begin{ghinhoanh}
\textbf{Short English to remember:}

Friendship lasts better when responsibility walks with it.

\end{ghinhoanh}

\begin{tuhocnhanh}
\textbf{Gợi ý để dạy và tự nghĩ / Teaching and reflection prompts}

1. Vì sao một chuyện tiền nhỏ lại làm rạn bạn bè? \textit{Why can a small money issue damage friendship?}

2. Bài học nào nên dạy học sinh ở đây? \textit{What lesson should students learn here?}

3. Ranh giới trong tình bạn cần bắt đầu từ đâu? \textit{Where should boundaries in friendship begin?}
\end{tuhocnhanh}
\ngancach

\section{Bị Cô Lập Giữa Một Lớp Đông}
\begin{truyen}{Bị Cô Lập Giữa Một Lớp Đông}{Feeling Isolated in a Full Classroom}
\chuhoa{C}{ó em ngồi giữa rất nhiều bạn nhưng gần như không có ai thật sự đứng cạnh.}
\textit{(Some students sit among many classmates yet have almost no one truly beside them.)}

Các em không bị đánh, không bị chửi công khai, chỉ bị bỏ ra ngoài mọi vòng trò chuyện.
\textit{(They are not openly attacked. They are simply left outside every circle of conversation.)}

Cái lạnh này rất khó thấy vì nó không ồn.
\textit{(This kind of coldness is hard to notice because it is quiet.)}

Nhưng nó làm tuổi đi học trở nên dài và nặng hơn rất nhiều.
\textit{(Yet it makes school life feel much longer and heavier.)}

\end{truyen}

\begin{baihoc}
Cô lập không cần tiếng động lớn mới làm người ta đau.
\textit{(Isolation does not need loud drama to hurt deeply.)}
\end{baihoc}

\begin{ghinhoanh}
\textbf{Short English to remember:}

Isolation does not need loud drama to hurt deeply.

\end{ghinhoanh}

\begin{tuhocnhanh}
\textbf{Gợi ý để dạy và tự nghĩ / Teaching and reflection prompts}

1. Vì sao kiểu cô lập này khó thấy? \textit{Why is this form of isolation hard to notice?}

2. Nó làm học sinh mệt ở điểm nào? \textit{How does it exhaust a student?}

3. Người dạy có thể can thiệp bằng cách nào? \textit{How can a teacher intervene?}
\end{tuhocnhanh}
\ngancach

\section{Chơi Chung Nhưng Không Thật Sự Thuộc Về}
\begin{truyen}{Chơi Chung Nhưng Không Thật Sự Thuộc Về}{Being in the Group Without Truly Belonging}
\chuhoa{C}{ó em giờ ra chơi vẫn đứng chung, vẫn cười theo, nhưng nhìn kỹ là biết em không thật sự ở trong nhóm.}
\textit{(Some students stand with the group at break time and even laugh along, but a closer look shows they do not truly belong there.)}

Tuổi học trò sợ nhất nhiều khi không phải bị ghét công khai.
\textit{(What students often fear most is not open hatred.)}

Mà là được cho đứng gần, nhưng không được thật sự bước vào.
\textit{(It is being allowed to stand near, but not truly invited in.)}

Sự lạc lõng kín đáo này làm mệt hơn người lớn tưởng.
\textit{(This quiet exclusion can exhaust students more than adults imagine.)}

\end{truyen}

\begin{baihoc}
Không phải cứ có người xung quanh là một đứa trẻ đã hết cô đơn.
\textit{(Being surrounded by people does not mean a child is no longer lonely.)}
\end{baihoc}

\begin{ghinhoanh}
\textbf{Short English to remember:}

Being surrounded by people does not always mean belonging.

\end{ghinhoanh}

\begin{tuhocnhanh}
\textbf{Gợi ý để dạy và tự nghĩ / Teaching and reflection prompts}

1. Dấu hiệu “ở gần mà chưa thuộc về” là gì? \textit{What are the signs of being near but not included?}

2. Vì sao kiểu cô đơn này khó nhận ra? \textit{Why is this kind of loneliness hard to notice?}

3. Người dạy có thể tạo chỗ thuộc về ra sao? \textit{How can a teacher create real belonging?}
\end{tuhocnhanh}
\ngancach

\section{Một Câu Trêu Có Thể Theo Em Cả Tuần}
\begin{truyen}{Một Câu Trêu Có Thể Theo Em Cả Tuần}{One Teasing Sentence Can Follow a Student for Days}
\chuhoa{C}{ó những câu bạn bè nói ra như đùa, nhưng người nhận không đùa nổi.}
\textit{(Some comments are said as jokes, but the person receiving them cannot laugh.)}

Một câu về ngoại hình, nhà cửa, điểm số, hay cách nói năng có thể theo một học sinh rất lâu.
\textit{(A comment about looks, family, grades, or speech can stay with a student for a long time.)}

Người nói quên nhanh. Người nghe không quên nhanh như vậy.
\textit{(The speaker forgets quickly. The listener does not.)}

Tuổi này da mỏng hơn người lớn tưởng.
\textit{(At this age, the skin of the heart is thinner than adults think.)}

\end{truyen}

\begin{baihoc}
Một lời trêu nhỏ với người này có thể là một vết day dứt với người khác.
\textit{(A small joke to one person can become a lasting sting to another.)}
\end{baihoc}

\begin{ghinhoanh}
\textbf{Short English to remember:}

A small joke to one person can become a lasting sting to another.

\end{ghinhoanh}

\begin{tuhocnhanh}
\textbf{Gợi ý để dạy và tự nghĩ / Teaching and reflection prompts}

1. Vì sao câu trêu có thể đi xa hơn người nói nghĩ? \textit{Why can teasing go further than the speaker imagines?}

2. Học sinh cần được dạy điều gì về lời nói? \textit{What should students learn about words?}

3. Mình có đang xem nhẹ những câu đùa nhỏ không? \textit{Are we underestimating small jokes?}
\end{tuhocnhanh}
\ngancach

\section{Nghỉ Chơi Một Mình Không Hẳn Là Thích Một Mình}
\begin{truyen}{Nghỉ Chơi Một Mình Không Hẳn Là Thích Một Mình}{Being Alone at Break Does Not Always Mean Preferring It}
\chuhoa{C}{ó em cứ nghỉ chơi là ngồi yên. Người lớn nhìn nhanh thì bảo chắc em hướng nội, chắc em thích yên.}
\textit{(Some students sit alone every break. Adults may quickly label them introverted or quiet by preference.)}

Nhưng có em không phải thích một mình. Chỉ là không biết chen vào đâu.
\textit{(But some do not prefer being alone. They simply do not know where they can enter.)}

Sự rụt rè của trẻ đôi khi cần một cây cầu nhỏ, chứ không cần bị gọi tên lớn.
\textit{(A child’s shyness may need a small bridge, not a loud label.)}

Nếu không có cây cầu ấy, em sẽ quen dần với việc đứng ngoài.
\textit{(Without that bridge, the child slowly gets used to staying outside.)}

\end{truyen}

\begin{baihoc}
Có những học sinh không cần ép hòa đồng. Các em cần một lối vào an toàn.
\textit{(Some students do not need forced sociability. They need a safe entry point.)}
\end{baihoc}

\begin{ghinhoanh}
\textbf{Short English to remember:}

Some students do not need to be pushed.

They need a safe way in.

\end{ghinhoanh}

\begin{tuhocnhanh}
\textbf{Gợi ý để dạy và tự nghĩ / Teaching and reflection prompts}

1. Em này cần gì hơn là lời gắn nhãn? \textit{What does this student need more than a label?}

2. “Lối vào an toàn” có thể là gì? \textit{What could a safe entry point look like?}

3. Mình có đang để em quen với việc đứng ngoài không? \textit{Are we letting the student get used to exclusion?}
\end{tuhocnhanh}
\ngancach

\section{Bạn Tốt Không Phải Lúc Nào Cũng Nói Lời Hay}
\begin{truyen}{Bạn Tốt Không Phải Lúc Nào Cũng Nói Lời Hay}{Good Friends Do Not Always Speak Sweetly}
\chuhoa{T}{uổi học trò có những tình bạn rất thật. Có em nhắc bạn học, cản bạn làm bậy, kéo bạn về khi bạn đang đi lệch.}
\textit{(School-age friendship can be very real. Some students remind their friends to study, stop them from doing wrong, or pull them back when they drift.)}

Lúc đó, lời nói nghe không ngọt.
\textit{(At those moments, the words do not sound sweet.)}

Nhưng không phải ai làm mình khó chịu cũng là người xấu.
\textit{(Not everyone who makes us uncomfortable is harmful.)}

Có những người bạn quý vì dám giữ mình lại đúng lúc.
\textit{(Some friends are precious because they dare to hold us back at the right moment.)}

\end{truyen}

\begin{baihoc}
Tình bạn tốt không chỉ ở chỗ cùng vui, mà còn ở chỗ cùng giữ nhau khỏi lệch đường.
\textit{(Good friendship is not only about sharing joy, but also about keeping each other from going astray.)}
\end{baihoc}

\begin{ghinhoanh}
\textbf{Short English to remember:}

Good friendship also means keeping each other from going astray.

\end{ghinhoanh}

\begin{tuhocnhanh}
\textbf{Gợi ý để dạy và tự nghĩ / Teaching and reflection prompts}

1. Dấu hiệu của một người bạn tốt là gì ngoài sự vui vẻ? \textit{What marks a good friend beyond shared fun?}

2. Vì sao lời khó nghe đôi khi lại đáng quý? \textit{Why can uncomfortable words be valuable?}

3. Mình có dạy học sinh phân biệt điều này không? \textit{Do we teach students how to tell the difference?}
\end{tuhocnhanh}
